% !TeX root = ../MAgent_Theory_2.0_companion.tex
\chapter{Hidden Conditional Information Under Smooth Dynamics}
\label{comp:smooth}

This chapter concerns the reciprocal law potential and its Fisher flow. An exact reduction need not identify agents with one another. The requirement is that discarded distinctions have no effect on the future quantities the reduced description promises to reproduce. Under smooth KL dynamics this requirement has a concrete content: hidden conditional information can contribute to the force on a retained marginal. A per-agent reduced state must retain that contribution, or have a reason why it cancels. The first part of the chapter builds such a state completion for a single agent's private information; parents, in the sense of grouped agents, appear only in the erasure example at the end.

\section{A hidden conditional law can move a visible marginal}

Write a law as $q(dz,dv)=\mu(dz)\kappa_z(dv)$. Here $z$ is retained and $v$ is internal. Against a reference $\nu(dz)A(dv)$, the KL chain rule is
\begin{equation}
 \KL(q\Vert\nu\otimes A)
 =\KL(\mu\Vert\nu)+\int H(z)\mu(dz),\qquad
 H(z)=\KL(\kappa_z\Vert A).
 \label{comp:smooth-energy}
\end{equation}
If one public value has a more costly hidden conditional law, changing its probability changes the full energy. The public Fisher velocity therefore contains $H(z)$ as well as the ordinary public force. A single average hidden entropy does not tell the reduced state which public values carry that cost.

There is an exact architecture in which this is the only hidden contribution. A sender transmits its public coordinate through a normalized, possibly law-dependent channel, and its private coordinate is not transmitted. The receiver reconstructs its private coordinate from its own conditional $\kappa_{i,z}$. The receiver law and compared source law then have the same private conditional, so that conditional cancels from their peer KL. With a fixed reference whose private factors are independent, the two-sector local model still couples the public belief and model laws. The original potential and Fisher metric give
\begin{equation}
 \begin{aligned}
 \dot\mu&=-\mu[h^c+H-\E_\mu(h^c+H)],\\
 \dot\ell_z(v)&=H(z)-\ell_z(v),\qquad
 \ell_z=\log(d\kappa_z/dA).
 \end{aligned}
 \label{comp:smooth-flow}
\end{equation}
The public score $h^c$ includes every incoming interaction and every channel derivative. Cancellation is a functional identity, so it does not require freezing a live channel while differentiating.

The second equation contains no public law and no other agent. Each private conditional simply relaxes toward its reference along an explicit path, whatever the rest of the population does, and the public flow feels it only as a time-dependent forcing fixed at the start. In this architecture private information never changes an interaction. That makes the construction solvable, and it is also its main limitation: the interesting case, in which private variables later influence interactions, is not covered.

\section{Retain the distribution of density values}

The conditional equation in \cref{comp:smooth-flow} moves a log-density value according to its current value and the conditional KL. It does not need the name of the private point carrying that value. This suggests retaining the probability law
\begin{equation}
 \xi_z=(\ell_z)_\#A,\qquad
 \int e^r\xi_z(dr)=1,\qquad
 H(z)=\int re^r\xi_z(dr).
 \label{comp:spectrum}
\end{equation}
For example, use a uniform reference on two private states. The laws $(3/4,1/4)$ and $(1/4,3/4)$ have the same spectrum: half the reference mass sees $\log(3/2)$ and half sees $\log(1/2)$. Their private preferences differ, but their contribution to the public force agrees.

What this forgets depends entirely on the declared reference. With the reference $(0.6,0.4)$ the same two laws have different spectra, and on a finite alphabet whose atoms have no coincidences among subset masses the spectrum identifies the conditional completely. The compression is the symmetry of the reference, not something the dynamics creates. In statistical language the spectrum is the law of the log-likelihood ratio between the conditional and the reference, so equal spectra mean that the two comparisons are equally informative in Blackwell's sense.

The complete spectrum evolves exactly. With $a=e^{-t}$,
\begin{equation}
 \xi_z(t)=\left(r\mapsto ar-\log\int e^{ar'}\xi_z(0,dr')\right)_\#\xi_z(0).
 \label{comp:spectrum-evolution}
\end{equation}
This follows by solving $d\kappa_z(t)/dA\propto(d\kappa_z(0)/dA)^a$ and normalizing. The formula applies to arbitrary private conditionals whose log densities are bounded uniformly in the public coordinate; on general public spaces the public flow also needs bounded log densities for its reference and channels whose forces are Lipschitz in those charts. It does not fit a Gaussian, and the retained state can remain infinite-dimensional. When the initial conditionals belong to an exponential family based at the reference, however, only their natural parameters change, shrinking by the factor $a$, so the private part then closes on finitely many numbers.

If all public values within an agent have the same conditional spectrum, $H(z)$ is constant and disappears under Fisher centering. The public laws then follow the original coarse KL dynamics exactly. The private conditionals can still differ: a public bit can select between the two swapped binary laws above. This gives internal correlation without changing the public force. Other agents may have different public laws and different spectra; equality between agents is unnecessary. For the public flow alone even less is needed: only how the R\'enyi divergences of the conditionals from the reference vary with $z$, up to a part common to all $z$.

This property is not created at a finite time by the conditional relaxation. At every finite time $a>0$, the power transformation and its normalization can be inverted. Different spectra therefore remain different until an asymptotic limit, unless some additional operation changes them. Retaining unequal spectra is an exact alternative, but their existence alone does not select groups of agents.

The architecture matters. If the channel carried the sender's private coordinate, two agents with the same spectra could sit at different peer costs: with a uniform two-point reference, the laws $(3/4,1/4)$ and $(1/4,3/4)$ compared through an identity private channel cost $\tfrac12\log3$, while two copies of the first cost nothing. Separate spectra per agent then no longer determine the forces.

\section{What a future record is allowed to inspect}

The same two swapped private laws show why a boundary protocol matters. Let record $1$ have probabilities $(3/4,1/4)$ at the two private points. The record probabilities are respectively $5/8$ and $3/8$. Equal unmarked spectra therefore do not preserve this response.

If all permitted private readouts use a fixed mark $T(v)$, retain the joint reference distribution of $(\ell_z(v),T(v))$. A likelihood $L(o\mid z,T(v))$ then gives the conditional record probability
\begin{equation}
 N_z(o)=\int e^r L(o\mid z,d)\eta_z(dr,dd),\qquad
 \eta_z=(\ell_z,T)_\#A.
\end{equation}
After the record, the log-density coordinate changes by
\begin{equation}
 r^+=r+\log L(o\mid z,d)-\log N_z(o).
\end{equation}
The public marginal is multiplied by $N_z(o)$ and normalized. Between records the same smooth conditional relaxation resumes, so the record's imprint on the private sector fades: the potential has no record term, and records act only at the moments they arrive. This gives an exact inherited record-and-update sequence for the declared protocols. If the marks identify every private point, however, the supposed compression may retain everything. Sufficiency is relative to what the future interaction can ask.

\section{A smooth channel can erase a dependence in finite time}

A separate four-agent example shows that exact independence can become available during the original smooth flow. Each belief and model law has a public bit, a private bit, and an internal gate bit. A fixed normalized local evaluator couples the public belief and model means. The gate means relax toward a fixed reference value. The private part of a peer channel depends on the source's private mean multiplied by a smooth function $\alpha$ of the population's gate means.

The function $\alpha$ has an exactly flat zero region. Its derivative is included in the Fisher equations. The gate is in effect an internal clock: it relaxes on its own, perturbed only by a term that a small private gain keeps dominated, and it enters the flat region at a finite time and stays there. Once inside, every peer's private cost in a receiver row is the same for all sources. It cancels from the Gibbs source comparison. The reduced energy and Fisher metric then split into a public part and an independent private-and-gate part. Both the finite-time erasure and the change of preferred source follow from a declared ordering of the private gain, the public peer strength, and the gate scale; the example does not show that the KL dynamics produces either without that ordering.

The public states start in two invariant sign regions, with two agents in each. No labels are assigned to groups by the equations. Initially, private similarities can select crossed pairs, each containing one member from each public region. After private erasure, the preferred sources change: each agent prefers the other member of its public region, and that preference persists. All cross-pair source weights remain positive. The two members of a preferred pair are distinct at every finite time, but this heterogeneity is transient: for a fixed temperature and a small enough public peer strength, their public states converge exponentially to each other, so each pair approaches a common state.

The interaction is not what produces this sequence. If every peer cost is deleted, the same rule for choosing preferred sources reads the crossed pairs first and the regional pairs afterward, because the private advantage and the public sign regions are supplied by the initial states; only the moment of the switch moves, since the gates then close more slowly. The example is therefore a single-stage selection with exact inheritance, not a formation.

For each pair, retain the empirical measure of its two paired public belief and model means. The original public vector field gives its exact evolution, and fresh public responses are computed by selecting and sampling a child from that measure. The private and gate laws can be forgotten after erasure, because they form a factor of the dynamics that has decoupled and that no admitted response reads. This is a noninjective reduction, but the empirical measure keeps every public child state, up to its label, and it is not itself a law pair of the children's type.

The example is an existence construction with an engineered channel, a product preparation on finite alphabets, and a specified class of boundary responses. It adds no fusion penalty and imposes no synchronization. It does not show that arbitrary smooth channels acquire a sufficient interface, that the distinct public child states admit a smaller aggregate description, or that the two resulting groups form another parent. Those are the remaining formation obligations.
